% §3 The Two-Floor Framework
We formalise responsiveness as two floors, model the orchestration that determines them, and
define the vocabulary used throughout the analysis.
All latency values refer to a fixed in-flight policy: a fixed endpointing configuration, a
fixed TTS buffer depth, and a fixed decode speed.

\begin{definition}[Responsiveness floors]
\label{def:floors}
For a single agent turn, \emph{response onset} is the interval from the detectable end of the
user's turn to the first audible unit of the agent's response.
\emph{Response cessation} is the interval from the first detectable onset of a user
interruption to the moment no further audibility of the agent's prior turn remains.
\end{definition}

Onset is the object of the companion Latency Floor Theorem and its Cascade Competitiveness
Criterion.
Cessation is the object of this paper.
The two intervals are disjoint in time but not in their dependencies: both read the same
detection subsystem, buffer policy, and decode speed.

\begin{definition}[Three clocks]
\label{def:clocks}
An orchestrator maintains three clocks over a turn.
The \emph{generation clock} tracks the language model's streamed output.
The \emph{delivery clock} tracks what has been handed to the speech synthesis and playback
path.
The \emph{tool clock} tracks tool execution scheduled in the background.
Control flow is expressed as events: a user turn, a tool completion, a delivery-started,
delivery-completed, or delivery-interrupted.
\end{definition}

The decoupling of the three clocks is the orchestration property under study.
A \emph{serial} topology couples all three clocks inside one call stack; a
\emph{concurrent} topology schedules them independently.

\begin{definition}[Interruption phases]
\label{def:phases}
An interruption lands in one of three phases of the agent's turn:
\emph{phase-gen} while generation runs but nothing is yet audible,
\emph{phase-play} while audio is being delivered, and
\emph{phase-tool} while the agent waits on a background tool.
\end{definition}

The phase determines what ``stopping'' means and which floor is re-entered afterwards.

\begin{definition}[Tool typing]
\label{def:tool_typing}
A tool is of
\emph{type A} if its result is required before any response audio can be produced,
\emph{type B} if its result can be deferred past the streaming-TTS lead-in, and
\emph{type C} if its result is needed only for a later turn.
\end{definition}

Type-A tools gate first audio and therefore sit on the onset floor; types B and C admit
hiding behind delivery lead-in or background execution.

\begin{definition}[Interruptibility]
\label{def:interruptibility}
An orchestrator is \emph{generation-interruptible} if cancelling the active model stream is a
first-class primitive.
It is \emph{execution-persistent} if cancelling generation does not cancel background tool
execution: pending tool calls continue, complete, and re-enter the model with their results.
\end{definition}

These two predicates jointly describe the capability that a cessation analysis operates on;
how \emph{fast} the capability is realised is set by the in-flight granularity of the
delivery and tool clocks.

\begin{definition}[API-valid history]
\label{def:apivalid}
A conversation history is \emph{API-valid} if it contains no assistant message whose tool-call
batch lacks its complete set of tool-result messages, for the model provider's function-calling
protocol.
\end{definition}

API-validity is the serializability constraint that tool persistence must respect: a pending
batch cannot be flushed to the provider, and cannot be abandoned without making the history
unplayable.

\subsection{An event system as the mechanism}
We instantiate the three clocks as an event-driven orchestrator: user texts and tool
completions are generic events; delivery status is reported back by the transport
(started, completed, interrupted); and interruption cancels the generation task without
touching the tool dispatcher.
Interrupted delivery is recorded and exposed as internal context on the next model
invocation, so the agent knows its previous response may not have been heard.
This instantiation is the object of the PoC validation in
Section~\ref{sec:analysis}; it is the configuration against which the serial topologies of
Section~\ref{sec:practical} are compared, not a uniqueness claim.
